% !TeX root = ../main.tex
\section{Cube Complex}
\label{sec:L08-W03D2}


\begin{definition}
    \label{def:cube-complex}
    A cube complex is a space built out of Euclidean \(k\)-cubes \([0,1]^k\) (\(k\) may be different) by gluing faces by (linear) isometries.
\end{definition}

\begin{example}[Examples of cube complex]
    \label{ex:cube-complex}
    \begin{itemize}
        \item \(\mathbb{S}^1\).
        \item Any \(1\)-dimensional CW complex.
        \item Tori.
        \item \(\mathbb{S}^2\)
        \item Closed orientable genus 2 surface \(\Sigma_2\) built out of \(8\) pentagons (glue 4 pentagons at each vertex). Note that by connecting the face center and five edge centers, we can subdivide a pentagon into 5 squares.
        \item Poincare homology sphere (See \cref{ex:Poincare-homology-sphere}).
        \item Seifert–Weber dodecahedral space (See \cref{rmk:Seifert-Weber-space}).
    \end{itemize}
\end{example}

\begin{definition}
    \label{def:link}
    In a cube complex \(X\), let \(v\) be a vertex. \(lk(v)\) is the boundary of an \(\epsilon\)-neighborhood of \(v\) in \(X\) (where \(epsilon\) is small, e.g. \(1/3\)). It consists of a collection of spherical simplices and is a \(\Delta\)-complex.
\end{definition}

\begin{example}
    \label{ex:link-of-cube-complex}
    \begin{itemize}
        \item For a \(1\)-simplex, \(lk(v)\) is a point (that is \(\epsilon\)-away from \(v\)).
        \item For a \(2\)-simplex, \(lk(v)\) is an arc (of a circle with radius \(\epsilon\)).
        \item Glue the two adjacent edges of \(v\) together, then \(lk(v)\) is a loop.
    \end{itemize}
\end{example}

\begin{definition}[Gromov's Link Condition]
    \label{def:NPC-cube-complex}
    A cube complex is \emph{non-positively curved (NPC)} if for any vertex \(v\), \(lk(v)\) is simplicial and flag.
\end{definition}

\begin{definition}
    \label{def:CAT-0-cube-complex}
    An NPC cube complex is \emph{\CATz} if it's simply-connected.
\end{definition}

Note that this means any NPC cube complex admits a \CATz cube complex universal cover.

\begin{theorem}[Gromov]
    \label{thm:connected-NPC-cube-complex-aspherical}
    All \CATz cube complexes are contractible. Hence all connected NPC cube complexes are aspherical.
\end{theorem}

Note that the converse is not true. Consider 3 squares glue to one vertex (a corner of the shell of a cube). This space is contractible, but the link is an empty triangle, hence this space is not NPC.

Observe that all graphs are NPC, as their links are disjoint points.

\begin{proposition}
    \label{prop:NPC-criterion-of-2-dim-cube-complex}
    Let \(X\) be a 2-dimensional cube complex. Observe that  all links are graph. The space is NPC if and only if each link is simplicial and has no loops of length 4.
\end{proposition}

\begin{lemma}
    \label{lem:flag-simplicial-complex-NPC-Salvetti-complex}
    If \(L\) is a flag simplicial complex, then its Salvetti complex \(X_L\) (See \cref{def:Salvetti-complex}) is NPC.
\end{lemma}

Note that the link of the only vertex in \(X_L\) is not \(L\). We will discuss the relationship between the link and \(L\) in the future.

\begin{proof}
    We claim that each \(k\)-simplex in \(L\) corresponds to a join of \(k+1\) \(\mathbb{S}^0\) in the link. 

    To see this, recall that each vertex in \(L\) gives a circle in \(X_L\), hence an \(\mathbb{S}^0\) in the link. Suppose now the claim is true for \((k-1)\)-simplex and consider a \(k\)-simplex in \(L\). Notice that \(k\)-simplex is the join of one vertex to a \((k-1)\)-simplex. In the Salvetti complex, we have a \((k+1)\)-torus, which is the direct product of a \(k\)-torus and a circle. Thus the link is the join of \(\mathbb{S}^0\) and the link of \(k\)-torus.

    As a result, the link is simplicial.

    If the link is not flag, then there exists an empty \((k+1)\)-torus in the Salvetti complex, which corresponds to a shell of \(k\)-simplex in \(L\). So \(L\) is not flag when the link of Salvetti complex is not flag.

    TODO more detail?
\end{proof}

\begin{example}[Examples of fundamental groups of NPC cube complex] ~\\
    \begin{itemize}
        \item Free groups.
        \item Free abelian groups.
        \item Fundamental groups of closed surfaces.
        \item (Khan-Markovic, Bergeron-Wise) If \(M\) is a closed hyperbolic \(3\)-manifold, then there exists an NPC cube complex \(X\) such that \(\pi_1(X) = \pi_1(M)\). This theorem is a key ingredient of Agol's Virtual Haken theorem and Virtual Fibering theorem. 
    \end{itemize}
\end{example}

An advantage of cube complex over simplicial complex is the hyperplanes.

\begin{definition}
    \label{def:hyperplane-of-cube-complex}
    Given an \(n\)-dimensional cube \(\sigma\) in \(X\), a midcube of \(\sigma\) is an \(n-1\)-dimensional unit cube running through the barycenter of \(\sigma\) and parallel to one of the faces of \(\sigma\). 
    
    Two edges are parallel to each other if they are the opposite sides of one square. Let \([e]\) be the equivalent class of parallel edges.
    The hyperplane dual to \([e]\) is the collection of midcubes which intersect edges in \([e]\).
\end{definition}

Hyperplanes are naturally immersed in the cube complex. For \CATz cube complex, the hyperplanes are embedded. Thus hyperplanes allows us to investigate some codimension-1 behavior of the spaces.

% special vs CAT(0)?
% CAT(0) is special 
% most interesting case is compact special 
% universal cover of special is CAT(0) (every NPC has CAT(0) )


% L is a flag simplicial cplx then XL is NPC,,, how to see it's join; okayy n torus is n-1 torus product S1 so that gives join of Lk n-1 and S0

% free abelian groups

% introduction of RAAGs Ruth Charney

% italiano

% Bruno  mateli 